\chapter{The Tree as Stack Machine}

The Christmas Tree is not only a divinatory layout. It is a reversible
symbolic-computational object: a stack machine with twenty-two registers,
a program counter that moves through seven lifecycle steps, and a LIFO
serialization that lets its state be stored, restored, branched, and revised.

This chapter formalizes the Tree as an abstract machine. The formalization is
not a replacement for the reading practices described earlier; it is a grammar
that lets the reader treat the Tree as a programmable instrument. The same
formalism connects the Tree to the family's wider computational philosophy:
Forth's stack-based execution, khipu's cord-and-knot memory, and the
reversible-state containers described in the programming chapter.

\section{The abstract machine}

A Christmas Tree machine $M$ is a tuple:

\[
M = (S, P, C, R, \mathit{T\_STOP})
\]

where

\begin{description}
\item[$S$] is the stack, a LIFO sequence of twenty-two card tokens in a declared
      order;
\item[$P$] is the program counter, an integer from $0$ to $6$ representing the
      seven lifecycle steps (catalytic presence through terminal integration);
\item[$C$] is the crown register, holding the single token at position $S_0$;
\item[$R$] is the root register, holding the single token at position $S_{21}$;
\item[$\mathit{T\_STOP}$] is the halting predicate, a boolean that becomes true
      when the program counter has completed step $6$ and the root register has
      been evaluated.
\end{description}

The stack $S$ is ordered from crown ($S_0$) to root ($S_{21}$). The five
unfolding rows occupy positions $S_1$ through $S_{20}$ in row-major order:
row~1 (positions $1$--$2$), row~2 ($3$--$5$), row~3 ($6$--$9$), row~4
($10$--$14$), row~5 ($15$--$20$).

\section{Lifecycle steps and the program counter}

The program counter $P$ advances through seven states:

\begin{center}
\begin{tabular}{cll}
\toprule
$P$ & Name & Operation \\
\midrule
0 & Catalytic presence & Load crown into $C$; begin clock \\
1 & First unfolding & Execute row~1 ($S_1, S_2$) \\
2 & Second unfolding & Execute row~2 ($S_3, S_4, S_5$) \\
3 & Third unfolding & Execute row~3 ($S_6, S_7, S_8, S_9$) \\
4 & Fourth unfolding & Execute row~4 ($S_{10}, \dots, S_{14}$) \\
5 & Fifth unfolding & Execute row~5 ($S_{15}, \dots, S_{20}$) \\
6 & Terminal integration & Load root into $R$; evaluate $T\_STOP$ \\
\bottomrule
\end{tabular}
\end{center}

Execution is one-way: $P$ increments from $0$ to $6$ and does not reverse.
The root is not a step that produces a new row; it is the halting evaluation
that receives the completed state.

\section{The instruction set}

Each lifecycle step executes a small instruction set whose operands are the
card tokens in the current row or register. The instructions are interpretive
rather than arithmetic: they produce observations, not numeric results.

\begin{description}
\item[\texttt{LOAD} $n$] Place token $S_n$ into the active register.
\item[\texttt{EXEC} $r$] Execute row $r$ as a compound instruction: read the
      tokens left-to-right, derive adjacency relations, and produce a row-level
      observation.
\item[\texttt{PAIR} $i$ $j$] Form the adjacency relation between tokens at
      positions $i$ and $j$, producing a directed observation (left-to-right
      or right-to-left as declared).
\item[\texttt{FAMILY} $r$] Read row $r$ as a functional family: identify
      whether the tokens belong to a shared archetypal group (threshold,
      relation, motion, disturbance, method, procession, root) and produce a
      compound observation.
\item[\texttt{HALT}] Evaluate the root register, set $T\_STOP$, and terminate
      the execution trace.
\end{description}

These instructions do not modify the stack. They are read-only operations that
produce a trace of observations. The stack itself is the program; the
instructions are the reader's attention moving through the program.

\section{Serialization: the reversible stack}

The Tree's entire state can be serialized into a single ordered sequence:

\begin{verbatim}
serialize(M) = [C, row_1, row_2, row_3, row_4, row_5, R]
             = [S_0, S_1, S_2, ..., S_20, S_21]
\end{verbatim}

This is a LIFO-compatible encoding: if the serialized sequence is stored as a
stack with $S_0$ at the bottom and $S_{21}$ at the top, restoring the Tree is
a matter of dealing from the top into the root, then row~5, and so on, until
the crown is placed last. In practice, the reader more often stores the
sequence crown-first and rebuilds the layout outward, but the mathematical
correspondence to a push-pop stack is exact.

Reversibility means:

\begin{enumerate}
\item Any Tree can be serialized without loss of positional information.
\item A serialized Tree can be restored to its exact layout.
\item A restored Tree can be branched: one or more tokens may be replaced while
      the positional architecture remains intact.
\item A branched Tree can be re-serialized and stored as a new program.
\end{enumerate}

This is the same reversibility principle that governs the card-set programming
described in the programming chapter. The Christmas Tree is simply a larger
and more structured program than a three-card line or a five-card cross.

\section{Forth and khipu: two neighboring grammars}

The Tree's stack-machine behavior is close to Forth, the concatenative language
in which operations consume and produce stack values without named variables. In
Forth, a program is a sequence of words that manipulate a parameter stack.
In the Tree, a layout is a sequence of tokens that manipulate an attention
stack: each row consumes the unresolved observations from the rows above and
produces new observations that flow toward the root.

The connection is not metaphorical. A Forth word can be defined that executes
a Christmas Tree: it pushes the crown, unfolds each row as a nested word, and
halts at the root. The formal definition might read:

\begin{verbatim}
: christmas-tree
  crown execute
  row-1 execute
  row-2 execute
  row-3 execute
  row-4 execute
  row-5 execute
  root halt ;
\end{verbatim}

The Tree is also close to khipu, the Andean cord-and-knot recording system in
which a primary cord carries pendant cords, and knots on those pendants encode
quantitative and qualitative information. In the Tree:

\begin{itemize}
\item The crown-to-root axis is the primary cord.
\item Each row is a pendant cord group: row~1 has two pendants, row~2 has three,
      and so on.
\item Each card is a knot: a token tied at a specific position that encodes both
      its archetypal value (the knot's type) and its positional meaning (the
      pendant's place in the group).
\item The entire Tree is a khipu program that can be rolled, stored, and
      unrolled without altering the knot sequence.
\end{itemize}

These connections are offered as neighboring lenses, not as claims that the Tree
is ``really'' Forth or ``really'' khipu. The native grammar of the Tree remains
card-based and lifecycle-based. The formal machine simply shows that the Tree's
operations are computationally coherent: they can be implemented, stored,
replayed, and reasoned about with the same rigor as any other stack machine.

\section{Execution trace: the realized Tree}

The realized Tree from Chapter~5 can be read as an execution trace. The
following table shows the program counter advancing through the lifecycle,
with the instruction executed at each step and the observation produced.

\begin{center}
\begin{longtable}{clp{8cm}}
\toprule
$P$ & Instruction & Observation \\
\midrule
\endhead
0 & \texttt{LOAD} $S_0$ & XIII Death at crown: catalytic release of an ended form \\
1 & \texttt{EXEC} row~1 & VI Lovers + VIII Justice: attraction becomes reciprocal choice under law \\
2 & \texttt{EXEC} row~2 & VII Chariot + XVIII Moon + XXI World: directed motion through obscurity toward integration \\
3 & \texttt{FAMILY} row~3 & XVI Tower + IV Emperor + XVII Star + III Empress: ruin cleared, rebuilt under authority, illuminated by hope, made fertile \\
4 & \texttt{EXEC} row~4 & I Magician + XII Hanged Man + V Pope + II Papess + XIV Temperance: craft suspended, transmitted, hidden, then blended \\
5 & \texttt{EXEC} row~5 & IX Hermit + X Wheel + 0 Fool + XI Strength + XX Judgement + XIX Sun: private light turns through chance, leaps beyond map, gathers power, receives call, becomes visible warmth \\
6 & \texttt{HALT} & XV Devil at root: attachment, appetite, and the terms of desire made conscious \\
\bottomrule
\end{longtable}
\end{center}

The trace is read-only: no instruction modifies the stack. But the trace itself
is a new artifact. It is the derivation that the reader carries back to the
parent question. The stack machine has produced a proof: given these twenty-two
tokens in these twenty-two positions, the lifecycle executes from catalytic
release to conscious attachment.

\section{Branching and revision}

A stored Tree can be revised without being redrawn from nothing. The reader may:

\begin{enumerate}
\item \textbf{Replace a parameter.} In a full-deck elaboration, a Marseille Minor
      beneath a Major position may be replaced while the Major architecture
      remains unchanged.
\item \textbf{Swap a Major.} The crown may be changed from XIII Death to 0 Fool,
      producing a different catalytic instruction while the row geometry stays
      intact.
\item \textbf{Insert an inspection.} A row may be expanded: one token becomes a
      parent, and a child English line or Marseille Minor is drawn, read, and
      returned before the main program resumes.
\item \textbf{Store multiple versions.} The same question may be programmed into
      several Trees, each with a different crown or root, and the versions may
      be compared after execution.
\end{enumerate}

Each branch is itself a reversible program. The reader serializes the original,
modifies one or more positions, serializes the branch, and records the
modification as a lineage note: ``branched from baseline Tree at $S_0$, crown
replaced, row~1 re-executed.''

\section{The stack as memory container}

A serialized Tree is a small memory container: twenty-two tokens in a declared
order, with a position map that reconstructs the geometry. In the family's
wider work on computational memory and continuity-bearing systems, a container
of this size is a \emph{cell}: a bounded body with local state, a birth
certificate (the question and date that produced it), and a conservation barrier
(the rule that the stack must not be altered during execution).

The Tree-as-cell has:

\begin{itemize}
\item \textbf{Identity:} the question that produced it;
\item \textbf{State:} the ordered stack and the program counter;
\item \textbf{History:} parent readings, expansions, and lineage notes;
\item \textbf{Lifecycle:} creation, execution, halting, storage, possible
      revision, and eventual composting (the deck is shuffled, the program
      returns to entropy).
\end{itemize}

When the Tree reaches $T\_STOP$, it does not die. It becomes soil: the
observations are recorded, the stack may be preserved or composted, and the
derived reading enters the reader's accumulated garden of tested patterns.

\section{The clock and the phase-space}

The Tree's execution is also a trace through a bounded phase-space. The
phase-space has twenty-two dimensions (one per position), but the program
counter constrains the path: the reader does not read all positions
simultaneously. The path is sequential, and the observations at each step
become constraints that narrow the possible meanings of later steps.

This is why the Tree produces detail through convergence rather than
accumulation. The phase-space is large, but the lifecycle steps are few and
directed. The reader's attention is the walker, and the stack is the terrain.

The formal notation for this trace is simple:

\[
\tau = [(P_0, O_0), (P_1, O_1), \dots, (P_6, O_6)]
\]

where each pair $(P_i, O_i)$ is a program-counter state and the observation
produced at that state. The full trace $\tau$ is the reading's derivation log.
It can be stored, compared with other traces, and audited for consistency:
does the row~3 family observation contradict the row~2 observation? If so, the
contradiction is itself information: the Tree contains a tension that a
simpler layout might have hidden.

\section{From formal machine to living practice}

The stack machine is a grammar, not a ritual. The reader need not think in
terms of registers and program counters while laying the cards. The formalism
becomes useful when:

\begin{itemize}
\item the reader wants to store a Tree and restore it exactly;
\item the reader wants to compare two Trees that address the same question;
\item the reader wants to explain the Tree's structure to another practitioner
      in precise terms;
\item the reader wants to implement the Tree in software, as a playable
      reference engine or a digital garden;
\item the reader wants to connect the Tree to other stack-based or cord-based
      systems in the family's wider work.
\end{itemize}

The living practice remains the same: shuffle, invite, lay the crown, unfold the
rows, meet the root, and ask what the whole life has become. The machine simply
shows that the practice is not vague. It is a discoverable grammar with exact
operations, reversible states, and a halting condition that the reader learns
to recognize by feel.

\begin{quote}
A seed plus lawful unfolding produces a world; a world that has reached its
whole may rest. The machine records the law; the reader lives the rest.
\end{quote}
